% Short running form; the full title is too wide for the head's right box.
\runninghead{The Word That Heals from Afar}
\properlane{word-at-a-distance}{The Word That Heals from Afar: Faith Taught to Rest on Christ's Word}
\label{sec:word-at-a-distance}

Heard from its centre, the Gospel turns on one exchange. A father asks Christ to
come down before his son dies. Christ answers with a rebuke to a faith that
waits on signs, and then gives him, in place of his presence, a word: \latin{Vade,
filius tuus vivit}. The father believes the word and goes, and on the road he
finds it true at the very hour it was spoken. St Gregory the Great and St John
Chrysostom read the exchange as Christ teaching faith to rest on his word.
Gregory finds the reason in the one who speaks: the father sought the bodily
presence of one who is absent nowhere, and the Lord healed by command what he
had made by his will. Chrysostom finds it in what the word does: it heals the
father's mind together with the son's body, so that he learns to attend to
Christ's teaching rather than to his miracles. Read with that Gospel, the whole
formulary becomes a school of the same faith. The Introit confesses that God's
judgments are true before it asks his mercy; the Collect asks a mind at rest;
the Epistle asks understanding of God's will, and the Alleluia sings a heart
made ready to accept it; the Gradual and the Communion
pray as men who hope in a word and wait for its due time; the Secret asks the
healing of mysteries in which Christ acts unseen; and the Postcommunion asks the
obedience that the Introit confessed it had not given. The Mass places those
who pray it where the father stood: asking for life, given a word, and sent
away to find it true.

\subsection{Presence sought, majesty forgotten: St Gregory the Great}

Gregory begins his homily by saying that the Gospel needs exhortation more than
exposition, and that only one thing in it needs explaining: why a man who came
to ask his son's salvation should hear, ``Unless you see signs and wonders, you
believe not.'' The man plainly believed, for \latin{neque enim ab eo quaereret
salutem quem non crederet salvatorem}: he would not have asked salvation of one
he did not believe a saviour. Gregory's answer is that the request itself shows
where the faith failed. \latin{Poposcit namque ut descenderet, et sanaret filium
ejus. Corporalem ergo praesentiam Domini quaerebat, qui per spiritum nusquam
deerat.} He asked the Lord to come down; so he was seeking the bodily presence
of one who, \latin{per spiritum}, was nowhere absent. \latin{Si enim perfecte
credidisset, procul dubio sciret quia non esset locus ubi non esset Deus}: had
he believed perfectly, he would have known that there is no place where God is
not. He gave honour, Gregory says, not to majesty but to bodily presence, and so
he believed the Lord able to heal and yet thought him absent from a dying son.
The Lord's answer corrects that faith by what it does: \latin{solo jussu salutem
reddidit qui voluntate omnia creavit}, he who created all things by his will
gave health by his command alone (\work{Homiliae in Evangelia} 28, 1).

Gregory's reason is theological, and it governs the reading. The word is enough
because of who speaks it. A healing by command is the creating will acting at a
distance that is no distance to God; the father is asked to believe not a
report about a cure but the speaker of the word. Gregory then turns the episode
against his hearers. The same Lord who would not go in body to the ruler's son
offered to go to the centurion's slave (Mt 8:7), and Gregory reads the contrast
as a rebuke to pride, which honours in men their wealth and rank and not the
nature in which they were made to God's image: \latin{Ecce de coelo venit qui
servo in terra occurrere non despicit}, he who came from heaven does not disdain
to go to meet a slave on earth (28, 2). He ends among the martyrs at whose tomb
he is preaching. Death and grief are everywhere, and the world that once held
men by delight is now so full of blows \latin{ut ipse nos jam mundus mittat ad
Deum}, that the world itself now sends us to God; therefore, \latin{Nolite \dots\
in vobismetipsis pensare quod habetis, sed quid estis}, weigh in yourselves not
what you have but what you are (28, 3).

\subsection{The father healed with the son: St John Chrysostom}

Chrysostom grants the ruler a faith, and finds it unsound. ``Yet the
very coming and beseeching Him was a mark of faith''; but ``if this ruler also
believed, yet he believed not entirely or soundly, as is clear from his
enquiring `at what hour the fever left him,' since he desired to know whether it
did so of its own accord, or at the bidding of Christ.'' The rebuke is therefore
aimed at the man's state of mind, and it is a kindness: ``It is for this that
Christ rebuketh him and toucheth his conscience, to show that His miracles were
wrought principally for the sake of the soul. For here He healeth the father,
sick in mind, no less than the son, in order to persuade us to give heed to Him,
not by reason of His miracles, but of His teaching. For miracles are not for the
faithful, but for the unbelieving and the grosser sort'' (\work{Homilies on St
John} 35, 2).

Chrysostom also explains why Christ offered to go to the centurion and would not
go to the ruler. ``Because in the former case faith had been perfected, and
therefore He undertook to go, that we might learn the rightmindedness of the
man; but here the nobleman was imperfect.'' The ruler ``knew not yet clearly
that even when absent He could heal'', and Christ, by not going, taught him
what the centurion had known of himself. His conclusion is the moral of the
reading: ``What now are we taught by these things? Not to wait for miracles,
nor to seek pledges of the Power of God.'' And he draws it into a household's
daily life. Many give thanks when a child recovers; they ought to give thanks
``even if they obtain it not'', since ``Whom the Lord loveth He chasteneth''
(35, 3).

Gregory and Chrysostom agree on the healing by word and differ on its reasons.
For Gregory, Christ declines the ruler's house to humble the pride that
honours rank; for Chrysostom, because the ruler's faith, unlike the
centurion's, was not yet perfected. Gregory teaches who it is that heals at a
distance; Chrysostom, what that healing does to the one who receives the word.
The reading needs both: the object of faith, and its growth.

\subsection{Faith that grows on the road: St Augustine, St Cyril and St Thomas}

The Fathers do not agree on how much the ruler believed when he came, and the
difference matters. St Augustine reads the rebuke at its full weight: Christ
``shows us a man lukewarm, or cold in faith, or of no faith at all'', and the
evangelist confirms it, since ``himself believed'' comes only after the hour is
known: ``it follows that when he was making the request he did not yet
believe'' (\work{Tractates on John} 16, 3). St Cyril of Alexandria sees a faith
real but childish: the ruler ``cometh as to One able to heal, but he
understandeth not yet that He is by Nature God''; ``Feeble indeed unto
understanding is the nobleman, for he is a child in his petition for grace''.
Yet Christ ``doth not reject our lack of apprehension; but benefiteth even the
stumbling'', and ``in Go thy way is Faith'': ``The one command of the Saviour
healeth two souls. For in the nobleman it worketh unwonted faith, the child it
rescueth from bodily death'' (\work{Commentary on John} II.5). Gregory, as
above, says that the ruler believed and doubted, and does not speak of his
faith's completion. Schuster, commenting on this Mass, says that the ruler ``had
faith in Jesus'', which Christ desired ``first of all to purify \dots\ from any
material element and from all idea of self-interest'', requiring him ``to
believe in his word without actually beholding the cure'' (p.~176).

St Thomas Aquinas gathers these answers into a sequence. The ruler's faith had
two defects: he did not believe in Christ's divine power, since God is present
everywhere (Jer 23:24) and could heal from afar, and, as Chrysostom says, he
doubted. Christ's \latin{Vade} is a command to dispose himself for grace by the
movement of his free will toward God. The evangelist's \latin{et ibat} means
that he went on growing in faith, because on the way of God not to go forward is
to go back. And Aquinas traces the stages: weak when he first interceded,
firmer when he called Jesus Lord, more perfect when he believed the word and
went, perfected when he knew the hour and believed with his house, as the path
of the just ``goeth forwards, and increaseth even to perfect day'' (Prov 4:18;
\work{Super Ioannem} c.~4, lect.~7). St Anthony of Padua, preaching on the same
Gospel, draws his lesson from the order of the two verbs: the believing comes
before the going, as the heart's assent comes before the deed.

On the ruler's first faith, then, Augustine and Gregory answer differently,
and Chrysostom, Cyril, Aquinas and Schuster stand between them. On what follows
the word, Augustine, Chrysostom, Cyril and Aquinas place a fuller faith after
it, and the Gospel's last clause bears them out: \latin{et credidit ipse, et
domus eius tota}. The reading's moral sense rests on that growth: faith is
given a word, acts on it, and is made whole by finding it true.

\subsection{The way, the will of God, and a heart made ready}

The Introit begins where the father's faith had to begin, with God and not
with the sign. Before it asks mercy it confesses that all God has done was done
\latin{in vero iudicio}, and St Augustine teaches that this is the order of the
saints' prayer: ``all His saints, when set in the midst of tribulation, have
first praised His righteousness, and so sought His blessings'' (\work{Enarrationes in Psalmos} 144, on v.~17). Its verse blesses those ``who walk in the law of the
Lord''. St Hilary of Poitiers reads the way of that verse as Christ, who said
\latin{Ego sum via}, and insists that it is not enough to enter it: it is
\latin{non solum ineunda \dots\ sed etiam pergenda}, a way to be walked to its
end, since no one on the road yet holds what it leads to (\work{Tractatus super
Psalmos}, in Ps.~118, \latin{Aleph} 2). Cassiodorus likewise takes the way to be
the Lord Saviour, walked not with the body's feet but with the steps of good
action (\work{Expositio Psalmorum} 118.1). Neither speaks of the Gospel; but
the father of the Gospel walks such a road. \latin{Credidit homo sermoni \dots\
et ibat}: he believed the word and set out on it, before he knew what he would
find at the end.

The Epistle names what such walking needs: \latin{intellegentes quae sit
voluntas Dei}, understanding what the will of God is. Aquinas makes God's will
the first principle of the prudence the Apostle asks, as the Lord's Prayer asks
\latin{Fiat voluntas tua} (\work{In Epistolam ad Ephesios} c.~5, lect.~6). The
father's going was such prudence: he acted on a will declared to him in a word.

The Alleluia sings the heart that takes such a will: \latin{Paratum cor meum,
Deus, paratum cor meum}. On the same words in Ps~56, St Augustine says that
``the patience of good men with preparation of heart accepteth the will of
God'', and he makes the readiness an answer to what others prepare: ``He hath
prepared pit to deceive, shall I not prepare heart to suffer? He hath prepared
pit to oppress, shall I not prepare heart to endure?'' (\work{Enarrationes in
Psalmos} 56, on v.~8). St Robert Bellarmine, quoted above, lets the readiness
reach to everything, ``ready to take anything cheerfully from your hand''. Neither
speaks of the Gospel, and the join is the editor's; but the father came with a
heart not yet ready in that sense. He had fixed the means of his son's healing,
that Christ should come down, and Christ gave him another: a word in place of a
presence. He took it and went, and on the road his heart was the one the
Alleluia sings, ready to take from God's hand what God gives. Chrysostom's
counsel that thanks are owed ``even if they obtain it not'' asks the same
readiness of every household. The verse ends \latin{gloria mea}, and Schuster
reads it of God, called the glory of the soul not only as the author of its
glory but ``also because he alone is the just judge of our merits'' (p.~176). Here too the join to the Gospel is the editor's: the honour
that Gregory says pride pays to rank, and that Christ would not pay to the
ruler's, is not the glory the Alleluia sings.

And the Collect asks for the mind he went home with, \latin{secura mente}, a
mind at rest before the proof arrives. It asks that mind through pardon, and
Cassiodorus, expounding the Alleluia's verse, says that the faithful say it
rightly when they trust that their sins are forgiven by his gift
(\work{Expositio Psalmorum} 107.2). On the editor's reading the ready heart and
the mind at rest are one disposition, and the Collect asks the gift that makes
it.

\subsection{Hope in a word, and medicine unseen}

The Gradual sets the eyes of all creatures on God, who gives them food
\latin{in tempore opportuno}. St Augustine reads the due time as a physician's
care: ``Just as when thou refreshest a sick man in due season, when he ought to
receive, then Thou givest''; and he adds for anyone who has asked and not yet
received, ``let his eyes wait for the food, which He giveth in due season \dots\
Thou delayest, not deniest'' (\work{Enarrationes in Psalmos} 144, on vv.~15--16). The
father waited a night, \latin{heri hora septima}, between the word and the news,
and the food came in its season.

The Communion is a prayer from the same posture. It asks God to remember his
word, \latin{in quo mihi spem dedisti}. St Augustine explains that God forgets
nothing: ``When therefore it is said unto Him, `O remember,' the desire of him
who prayeth is displayed, because he asketh for what was promised \dots\ that
is, fulfil Thy promise to Thy servant'' (\work{Enarrationes in Psalmos} 118, sermo 15, on v.~49). St Robert Bellarmine adds why the promise must be asked: God ``still
wishes his faithful servants to ask him to carry them out; and thus, prayer
becomes one of the means through which God decreed to fulfill his promises''.
The father's double petition was such a prayer, and it was answered with a
word to be believed. The verse the Communion leaves unsung completes the
thought, \latin{quia eloquium tuum vivificavit me}, because thy word has given me
life, and the Gospel's word is \latin{Filius tuus vivit}.

Between them the Offertory and the Secret carry the father's situation into the
Mass. The Offertory weeps for a Sion that is remembered and not seen; faith of
this kind lives on a word held in absence. The Secret asks that the mysteries
furnish heavenly medicine and purge the heart. Dom Lucien Fromage, in the
continuation of \work{The Liturgical Year}, joins it to this Gospel: ``The whole
power of the God, who, with a word, cures both soul and body, resides in the
Mysteries which are about to be celebrated on our Altar here''; and he draws the
reading's moral: ``our Jesus has no need to come down from heaven to earth, in
order to give efficiency to the commands of his gracious will'' (\work{The
Liturgical Year} XI, p.~451). Schuster, too, reads the Secret of the
Eucharist as medicine, and the Gradual of the bread God prepares ``in every
region of the earth'' under ``the veil of faith, in order that we may have all
the merit of believing in the pure and single word of the Son of God''
(pp.~175--177). The Postcommunion closes the school: \latin{fac nos \dots\ tuis
semper oboedire mandatis}. The word received becomes the word obeyed, and the
disobedience confessed at the Introit is asked away at the end.

\subsection{The rebuke and the signs: the strongest difficulty}

The rebuke seems hard. A father whose son is dying asks for help and is told,
``Unless you see signs and wonders, you believe not.'' Chrysostom, Cyril and
Aquinas answer that the rebuke served him: it ``toucheth his conscience''
(Chrysostom), Christ ``benefiteth even the stumbling'' (Cyril), and the
reproach falls on a man instructed in the law who sought to believe by signs
rather than by the authority of Scripture, since signs are given to unbelievers
and the faithful are led by the Scriptures (Aquinas, after 1 Cor 14:22).
Augustine does not soften it: the rebuke exposes the man and, beyond him, a
country that needed signs, and Augustine sets him beside Thomas, to whom it was
said, ``Blessed are they that have not seen, and yet have believed'' (16, 4).
The reading keeps the rebuke at its weight and keeps also what followed it: the
same Christ who rebuked the man gave him the word.

Nor does the reading make signs worthless. The healing is itself a sign, the
second at Cana. Gregory says elsewhere that visible miracles were given
\latin{ut corda videntium ad fidem invisibilium pertrahant}, to draw the hearts
of those who saw them to faith in what is unseen, and that \latin{linguae in
signum sunt non fidelibus, sed infidelibus} (\work{Homiliae in Evangelia} 4, 3);
Chrysostom, that miracles are ``for the unbelieving''. The claim of the reading
concerns what faith rests on, not whether God gives signs. The father was given
both, and the order matters: the word first, believed; the sign after, found
true at the hour of the word.

\subsection{The four senses of this reading}

\begin{fourSenses}
\item[Literal.] A royal official of Capharnaum comes to Jesus at Cana and asks
him to come down and heal his dying son. He is rebuked for a faith that waits on
signs, receives instead the word ``Thy son liveth'', believes it and goes; on
the way he learns that the fever left the boy at the hour the word was spoken,
and he believes with his whole house. The Fathers read the rebuke as addressed
to a faith that wanted presence and proof (Augustine; Gregory; Chrysostom).

\item[Allegorical.] Christ, absent nowhere and healing by command what he
created by his will (Gregory), heals by his word without coming down, as he now
acts in his Church through word and mysteries without being seen (Fromage, of
the Secret; Schuster, of the Gradual). The seventh hour at which the fever left
the boy signifies the sevenfold gift of the Holy Spirit, by whom sins are
forgiven (Aquinas).

\item[Moral.] Believe the word before you see, and set out on it; let faith grow
from need to trust to the faith of a whole house (Aquinas; Cyril); do not wait
for miracles, and give thanks whether a child is spared or not (Chrysostom);
take from God's hand whatever he gives, with the heart the Alleluia calls ready
(Augustine; Bellarmine); weigh not what you have but what you are, and honour
in every person the image of God rather than rank, as the Lord would go to a
slave and not to a ruler's son (Gregory).

\item[Anagogical.] The father's road home, between the word and the sight of
its fulfilment, is the life of faith, which ends when the promise is kept and
seen: the Communion asks God to fulfil his promise (Augustine), and ``the
blessed enjoy that same truth which sustains us here on earth; but they enjoy
it without any veil, face to face'' (Schuster, p.~176); and the glory the
Alleluia sings is God himself, author of the glory that will make the soul
happy for all eternity (Schuster, on the Alleluia). The word that said
\latin{vivit} points to the life that the word of God gives without end.
\end{fourSenses}
